\item \textbf{Retroceso del núcleo emisor de radiación gamma.}
Cuando un núcleo excitado vuelve a su estado fundamental emitiendo un fotón gamma, debe retroceder para conservar el momento lineal.
\begin{enumerate}
    \item Demuestre que la energía de retroceso del núcleo emisor está dada por:
    $$ E_r = \frac{(\Delta E)^2}{2 M c^2} $$
    donde $\Delta E$ es la diferencia de energía de transición del núcleo y $M$ es su masa (asuma $hf \ll M c^2$).
    \item Calcule la energía de retroceso para un núcleo de $^{57}\text{Fe}$ ($M \approx 57\text{ u}$) que decae desde su estado excitado de $14.4\text{ kB}$ por emisión gamma.
\end{enumerate}